% !TITLE 牡丹亭·惊梦（节选）
% !AUTHOR 汤显祖
% !DYNASTY 明清
% !GENRE 戏曲唱词
% !ORDER 2

\begin{poem}
原来姹紫嫣红\textsuperscript{1}开遍，似这般都付与断井颓垣\textsuperscript{2}。
良辰美景奈何天\textsuperscript{3}，赏心乐事谁家院。
朝飞暮卷\textsuperscript{4}，云霞翠轩\textsuperscript{5}；雨丝风片，烟波画船——锦屏人\textsuperscript{6}忒\textsuperscript{7}看的这韶光\textsuperscript{8}贱！

遍青山啼红了杜鹃\textsuperscript{9}，荼蘼\textsuperscript{10}外烟丝醉软。
牡丹虽好，他春归怎占的先\textsuperscript{11}！

不到园林，怎知春色如许\textsuperscript{12}？
\end{poem}

\begin{annotations}
\notetext{1}{姹紫嫣红：形容各色花朵艳丽绚烂。姹（chà），美丽；嫣（yān），娇艳。}
\notetext{2}{断井颓垣：残破的水井和倒塌的墙壁。指废弃的院落，与“姹紫嫣红”形成强烈对比。}
\notetext{3}{奈何天：无可奈何的天气（或时运）。良辰美景虽好，天却未必作美——一种“美好的东西都与我无关”的怅惘。}
\notetext{4}{朝飞暮卷：语出王勃《滕王阁序》“画栋朝飞南浦云，珠帘暮卷西山雨”。形容楼阁之美。}
\notetext{5}{翠轩：翠绿色的轩窗或亭阁。轩（xuān），有窗的长廊或小室。}
\notetext{6}{锦屏人：锦缎屏风后面的人。指深闺中的女子（杜丽娘自指），被封闭在绣阁中，看不见外面的春光。}
\notetext{7}{忒：太、过于。忒（tuī）。}
\notetext{8}{韶光：美好的春光，也指青春年华。}
\notetext{9}{啼红了杜鹃：杜鹃花（又名映山红）开遍青山，像被杜鹃鸟啼出的血染红。杜鹃鸟啼声凄切，传说啼血染红花朵。}
\notetext{10}{荼蘼：一种晚春开花的蔷薇科植物。荼蘼（tú mí），花白色，开在春末，古人说“开到荼蘼花事了”——荼蘼开完，春天就结束了。}
\notetext{11}{他春归怎占的先：牡丹开在暮春，等它开放时春天都快过去了，怎么能算占了春光之先呢？杜丽娘以牡丹自比青春被耽误。}
\notetext{12}{如许：如此、这般。}
\end{annotations}

\begin{appreciation}
《牡丹亭》是明代汤显祖的传奇，写官家小姐杜丽娘为一场梦里的爱情死去、又为它复活的传说。这段唱词是全剧最有名的一段，也是中国戏曲里写青春觉醒最动人的篇章之一。

杜丽娘是南安太守的女儿，十六岁，被关在深闺里，连自己家的后花园都没进去过。这一天她第一次走进园子，开口就是"原来姹紫嫣红开遍"——"原来"两个字，就是一场震动：花开了这么多年，她今天才看见。

"似这般都付与断井颓垣"——这么好的花，对着的是破井塌墙，没有人看。美和荒废并排摆在一起，她看见的哪里是花，分明是自己：十六岁的春光，是不是也要像这园子一样白白荒掉？"良辰美景奈何天，赏心乐事谁家院"——古人说良辰、美景、赏心、乐事四样很难凑齐，杜丽娘一样都没有。

"朝飞暮卷，云霞翠轩；雨丝风片，烟波画船"，园子其实美得很。可一句"锦屏人忒看的这韶光贱"——锦屏人就是屏风后面的人，也就是她自己：深闺里的女儿，太把春光看得贱了。这一句骂的是谁？骂的是把她关在屏风里的世道，也骂自己从前竟不觉得。

"遍青山啼红了杜鹃"，杜鹃花开遍青山，红得像杜鹃鸟啼出的血，美得凄厉。"荼蘼外烟丝醉软"，荼蘼开在春末——开到荼蘼花事了，春天要完了。"牡丹虽好，他春归怎占的先"——牡丹是花中之王，却开在暮春，春天都快过完了才开，再高贵也迟了。杜丽娘以牡丹自比：我的青春，是不是也被耽误了？

"不到园林，怎知春色如许"——不进园子，永远不知道春天是这样的。一句大白话，是全剧的题眼。想想《诗经》里的女子：《桃夭》的姑娘"桃之夭夭，灼灼其华"，出门出嫁，春天看得多么坦荡；《关雎》的君子淑女，爱得光明正大。到了明代，一个姑娘连看一眼春天都要偷偷摸摸，要爱一个人只能到梦里去——同一个春天，隔了两千年，变成了两种过法。汤显祖不信这套，他说"情不知所起，一往而深，生者可以死，死可以生"——礼教关得住人，关不住情。

你现在的春天，出门就是；杜丽娘的春天，要拿命去换。所以别嫌花开花落是寻常事——你替她多看看，她那一扇门一辈子没打开，你的门一直开着。
\end{appreciation}
